\chapter{2014年天津卷（理）}

\section{单选题}

\begin{problem}
\(\i\) 是虚数单位，复数 \(\frac{7 + \i}{3 + 4\i} =\)（\quad）

\choices
    {\(1 - \i\)}
    {\(-1 + \i\)}
    {\(\frac{17}{25} + \frac{31}{25}\i\)}
    {\(-\frac{17}{7} + \frac{25}{7}\i\)}
\end{problem}

\begin{answer}
A
\end{answer}

\begin{solution}
分子、分母同乘以分母的共轭复数，得
\[
\frac{7+\i}{3+4\i}=\frac{(7+\i)(3-4\i)}{(3+4\i)(3-4\i)}=\frac{25-25\i}{25}=1-\i.
\]
因此答案为 A.
\end{solution}

\begin{problem}
设变量 \(x\)，\(y\) 满足约束条件 \(\begin{cases}
x + y - 2 \geqslant 0,\\
x - y - 2 \leqslant 0,\\
y \geqslant 1,
\end{cases}\) 则目标函数 \(z = x + 2y\) 的最小值为（\quad）

\choices
    {\(2\)}
    {\(3\)}
    {\(4\)}
    {\(5\)}
\end{problem}

\begin{answer}
B
\end{answer}

\begin{solution}
约束条件等价于 \(x\geqslant2-y\)，\(x\leqslant y+2\)，\(y\geqslant1\). 对于固定的 \(y\)，目标函数 \(z=x+2y\) 随 \(x\) 增大而增大，所以取 \(x=2-y\) 时 \(z\) 最小. 此时 \(z=2-y+2y=2+y\geqslant3\). 当 \(y=1\)，\(x=1\) 时各约束均成立，且 \(z=3\)，故最小值为 \(3\)，答案为 B.
\end{solution}

\begin{problem}
阅读下面的程序框图，运行相应的程序，输出的 \(S\) 的值为（\quad）

\texfigure[max width=0.50\linewidth]{img_repaint/2014/tianjin_science/q3_fig1.tex}

\choices
    {\(15\)}
    {\(105\)}
    {\(245\)}
    {\(945\)}
\end{problem}

\begin{answer}
B
\end{answer}

\begin{solution}
初始时 \(S=1\)，\(i=1\). 第一次循环中 \(T=2\times1+1=3\)，所以 \(S=3\)，随后 \(i=2\)；第二次循环中 \(T=5\)，所以 \(S=15\)，随后 \(i=3\)；第三次循环中 \(T=7\)，所以 \(S=105\)，随后 \(i=4\). 此时满足 \(i\geqslant4\)，程序输出 \(S=105\)，答案为 B.
\end{solution}

\begin{problem}
函数 \( f(x) = \log_{\frac{1}{2}}(x^2 - 4) \) 的单调递增区间是（\quad）

\choices
    {\((0, +\infty)\)}
    {\((-\infty, 0)\)}
    {\((2, +\infty)\)}
    {\((-\infty, -2)\)}
\end{problem}

\begin{answer}
D
\end{answer}

\begin{solution}
由对数真数为正，得 \(x^2-4>0\)，所以定义域为 \((-\infty,-2)\cup(2,+\infty)\). 因为底数 \(\frac12\) 小于 \(1\)，函数 \(y=\log_{1/2}u\) 是减函数；而 \(u=x^2-4\) 在 \((-\infty,-2)\) 上减，在 \((2,+\infty)\) 上增. 因此复合函数在 \((-\infty,-2)\) 上增，在 \((2,+\infty)\) 上减，单调递增区间为 \((-\infty,-2)\)，答案为 D.
\end{solution}

\begin{problem}
已知双曲线 \(\frac{x^2}{a^2} - \frac{y^2}{b^2} = 1\)（\(a > 0\)，\(b > 0\)）的一条渐近线平行于直线 \(l: y = 2x + 10\)，双曲线的一个焦点在直线 \(l\) 上，则双曲线的方程为（\quad）

\choices
    {\(\frac{x^2}{5} - \frac{y^2}{20} = 1\)}
    {\(\frac{x^2}{20} - \frac{y^2}{5} = 1\)}
    {\(\frac{3x^2}{25} - \frac{3y^2}{100} = 1\)}
    {\(\frac{3x^2}{100} - \frac{3y^2}{25} = 1\)}
\end{problem}

\begin{answer}
A
\end{answer}

\begin{solution}
双曲线的渐近线斜率为 \(\pm\frac ba\). 由一条渐近线平行于 \(y=2x+10\)，得 \(\frac ba=2\)，即 \(b=2a\). 设半焦距为 \(c\)，则 \(c^2=a^2+b^2=5a^2\). 双曲线的焦点为 \((\pm c,0)\)，在直线 \(y=2x+10\) 上的焦点必须满足 \(0=2x+10\)，故该焦点为 \((-5,0)\)，从而 \(c=5\). 因此 \(5a^2=25\)，得 \(a^2=5\)，\(b^2=20\). 双曲线方程为 \(\frac{x^2}{5}-\frac{y^2}{20}=1\)，答案为 A.
\end{solution}

\begin{problem}
如图，\(\triangle ABC\) 是圆的内接三角形，\(\angle BAC\) 的平分线交圆于点 \(D\)，交 \(BC\) 于 \(E\)，过点 \(B\) 的圆的切线与 \(AD\) 的延长线交于点 \(F\)，在上述条件下，给出下列四个结论：

\bitmapfigure[max width=0.42\linewidth]{img/2014/tianjin_science/q6_fig1.png}

\begin{circlelist}
\item \(BD\) 平分 \(\angle CBF\)；
\item \(FB^{2} = FD \cdot FA\)；
\item \(AE \cdot CE = BE \cdot DE\)；
\item \(AF \cdot BD = AB \cdot BF\).
\end{circlelist}

则所有正确结论的序号是（\quad）

\choices
    {\circled{1}\circled{2}}
    {\circled{3}\circled{4}}
    {\circled{1}\circled{2}\circled{3}}
    {\circled{1}\circled{2}\circled{4}}
\end{problem}

\begin{answer}
D
\end{answer}

\begin{solution}
令 \(\alpha=\angle BAD=\angle DAC\).

\begin{enumerate}
\item 由圆周角定理，\(\angle CBD=\angle CAD=\alpha\). 由切线与弦所夹的角等于弦所对的圆周角，\(\angle DBF=\angle BAD=\alpha\). 因而 \(BD\) 平分 \(\angle CBF\).
\item 直线 \(FAD\) 是过点 \(F\) 的割线，且 \(FB\) 是切线. 由切割线定理，\(FB^2=FA\cdot FD\)，所以结论 \circled{2} 正确.
\item \(AD\) 与 \(BC\) 交于 \(E\)，由相交弦定理，实际有 \(AE\cdot DE=BE\cdot CE\)，并非题给的 \(AE\cdot CE=BE\cdot DE\). 例如在等边三角形的情形，\(E\) 是 \(BC\) 的中点，而 \(D\) 是弧 \(BC\) 的中点；若圆半径为 \(R\)，则 \(AE=\frac{3R}{2}\)，\(DE=\frac R2\)，且 \(BE=CE\)，故结论 \circled{3} 不成立.
\item 设 \(\delta=\angle ADB\). 由切线与弦所夹的角定理，\(\sin\angle ABF=\sin\delta\)，又因 \(A,D,F\) 共线，\(\angle BAF=\angle BAD=\alpha\). 在三角形 \(ABF\) 和 \(ABD\) 中分别使用正弦定理，得
\[
\frac{AF}{BF}=\frac{\sin\angle ABF}{\sin\angle BAF}=\frac{\sin\delta}{\sin\alpha}=\frac{AB}{BD}.
\]
因此 \(AF\cdot BD=AB\cdot BF\)，结论 \circled{4} 正确.
\end{enumerate}

正确结论为 \circled{1}\circled{2}\circled{4}，答案为 D.
\end{solution}

\begin{problem}
设 \(a\)，\(b\in\R\)，则“\(a > b\)”是“\(a|a| > b|b|\)”的（\quad）

\choices
    {充分不必要条件}
    {必要不充分条件}
    {充要条件}
    {既不充分也不必要条件}
\end{problem}

\begin{answer}
C
\end{answer}

\begin{solution}
函数 \(g(t)=t|t|\) 在 \([0,+\infty)\) 上为 \(t^2\)，在 \((-\infty,0)\) 上为 \(-t^2\)，两段都严格递增，且在 \(0\) 处连续，因此 \(g\) 在 \(\R\) 上严格递增. 所以 \(a>b\) 当且仅当 \(g(a)>g(b)\)，即 \(a|a|>b|b|\). 因而“\(a>b\)”是“\(a|a|>b|b|\)”的充要条件，答案为 C.
\end{solution}

\begin{problem}
已知菱形 \(ABCD\) 的边长为 \(2\)，\(\angle BAD = 120^{\circ}\)，点 \(E\)，\(F\) 分别在边 \(BC\)，\(DC\) 上，\(\overrightarrow{BE} = \lambda\overrightarrow{BC}\)，\(\overrightarrow{DF} = \mu\overrightarrow{DC}\)，若 \(\overrightarrow{AE} \cdot \overrightarrow{AF} = 1\)，\(\overrightarrow{CE} \cdot \overrightarrow{CF} = -\frac{2}{3}\)，则 \(\lambda + \mu =\)（\quad）

\choices
    {\(\frac{1}{2}\)}
    {\(\frac{2}{3}\)}
    {\(\frac{5}{6}\)}
    {\(\frac{7}{12}\)}
\end{problem}

\begin{answer}
C
\end{answer}

\begin{solution}
设 \(\bs u=\overrightarrow{AB}\)，\(\bs v=\overrightarrow{AD}\). 由菱形边长和夹角，\(\bs u\cdot\bs u=\bs v\cdot\bs v=4\)，\(\bs u\cdot\bs v=4\cos120^\circ=-2\).

由点 \(E,F\) 的位置关系，有 \(\overrightarrow{AE}=\bs u+\lambda\bs v\)，\(\overrightarrow{AF}=\bs v+\mu\bs u\)，\(\overrightarrow{CE}=(\lambda-1)\bs v\)，\(\overrightarrow{CF}=(\mu-1)\bs u\). 因而两个数量积条件分别给出
\[
4\lambda+4\mu-2\lambda\mu-2=1.
\]
\[
-2(\lambda-1)(\mu-1)=-\frac23.
\]
令 \(s=\lambda+\mu\)，\(p=\lambda\mu\)，上式化为 \(4s-2p=3\)，\(p-s=-\frac23\). 两式联立消去 \(p\)，得
\[
s=\frac56.
\]
所以 \(\lambda+\mu=\frac56\)，答案为 C.
\end{solution}

\section{填空题}

\begin{problem}
某大学为了解在校本科生对参加某项社会实践活动的意向，拟采用分层抽样的方法，从该校四个年级的本科生中抽取一个容量为 \(300\) 的样本进行调查. 已知该校一年级，二年级，三年级，四年级的本科生人数之比为 \(4:5:5:6\)，则应从一年级本科生中抽取\fillinblank{} 名学生.
\end{problem}

\begin{answer}
\(60\)
\end{answer}

\begin{solution}
四个年级本科生人数的总份数为 \(4+5+5+6=20\). 按分层抽样，各年级抽取人数与该年级人数成正比，因此一年级应抽取
\[
300\times\frac4{20}=60
\]
名学生.
\end{solution}

\begin{problem}
已知一个几何体的三视图如图所示（单位：\(\mathrm{m}\)），则该几何体的体积为\fillinblank{}\(\mathrm{m}^{3}\).

\bitmapfigure[max width=0.58\linewidth]{img/2014/tianjin_science/q10_fig1.png}
\end{problem}

\begin{answer}
\(\frac{20\pi}{3}\)
\end{answer}

\begin{solution}
由正视图和侧视图可知，几何体下部是底面半径为 \(1\)，高为 \(4\) 的圆柱，上部是底面半径为 \(2\)，高为 \(2\) 的圆锥；俯视图中的外圆和内圆分别对应半径 \(2\) 和半径 \(1\). 因此几何体的体积为
\[
V=\pi\cdot1^2\cdot4+\frac13\pi\cdot2^2\cdot2=4\pi+\frac{8\pi}{3}=\frac{20\pi}{3}\ \mathrm{m}^{3}.
\]
\end{solution}

\begin{problem}
设 \(\{a_{n}\}\) 是首项为 \(a_{1}\)，公差为 \(-1\) 的等差数列，\(S_{n}\) 为其前 \(n\) 项和. 若 \(S_{1}\)，\(S_{2}\)，\(S_{4}\) 成等比数列，则 \(a_{1}\) 的值为\fillinblank{}.
\end{problem}

\begin{answer}
\(-\frac{1}{2}\)
\end{answer}

\begin{solution}
由等差数列通项公式，\(a_n=a_1-(n-1)\). 所以
\(S_1=a_1\)，\(S_2=a_1+(a_1-1)=2a_1-1\)，\(S_4=a_1+(a_1-1)+(a_1-2)+(a_1-3)=4a_1-6\). 由于 \(S_1,S_2,S_4\) 成等比数列，有
\[
(2a_1-1)^2=a_1(4a_1-6).
\]
展开并化简得 \(4a_1^2-4a_1+1=4a_1^2-6a_1\)，所以 \(2a_1+1=0\)，即 \(a_1=-\frac12\).
\end{solution}

\begin{problem}
在 \(\triangle ABC\) 中，内角 \(A\)，\(B\)，\(C\) 所对的边分别是 \(a\)，\(b\)，\(c\)，已知 \(b - c = \frac{1}{4}a\)，\(2\sin B = 3\sin C\)，则 \(\cos A\) 的值为\fillinblank{}.
\end{problem}

\begin{answer}
\(-\frac{1}{4}\)
\end{answer}

\begin{solution}
由正弦定理，\(\frac bc=\frac{\sin B}{\sin C}\). 由 \(2\sin B=3\sin C\)，得 \(b:c=3:2\). 设 \(b=3k\)，\(c=2k\)，其中 \(k>0\). 由 \(b-c=\frac14a\)，得 \(k=\frac14a\)，所以 \(a=4k\). 于是
\[
\cos A=\frac{b^2+c^2-a^2}{2bc}=\frac{(3k)^2+(2k)^2-(4k)^2}{2\cdot3k\cdot2k}=-\frac14.
\]
\end{solution}

\begin{problem}
在以 \(O\) 为极点的极坐标系中，圆 \(\rho = 4 \sin \theta\) 和直线 \(\rho \sin \theta = a\) 相交于 \(A\)，\(B\) 两点. 若 \(\triangle AOB\) 是等边三角形，则 \(a\) 的值为\fillinblank{}.
\end{problem}

\begin{answer}
\(3\)
\end{answer}

\begin{solution}
令 \(x=\rho\cos\theta\)，\(y=\rho\sin\theta\). 圆 \(\rho=4\sin\theta\) 化为 \(x^2+y^2=4y\)，直线 \(\rho\sin\theta=a\) 化为 \(y=a\). 设两个交点为 \(A=(-t,a)\)，\(B=(t,a)\)，其中 \(t>0\). 由交点在圆上，得
\[
t^2+a^2=4a.
\]
三角形 \(AOB\) 为等边三角形，所以 \(OA^2=AB^2\)，即
\[
t^2+a^2=(2t)^2,
\]
从而 \(a^2=3t^2\). 代入前式，得 \(\frac43a^2=4a\)，所以 \(a=0\) 或 \(a=3\). 当 \(a=0\) 时 \(t=0\)，两个交点重合，不符合题意，故 \(a=3\).
\end{solution}

\begin{problem}
已知函数 \( f(x) = |x^2 + 3x| \)，\( x \in \R \). 若方程 \( f(x) - a|x - 1| = 0 \) 恰有 4 个互异的实数根，则实数 \( a \) 的取值范围为\fillinblank{}.
\end{problem}

\begin{answer}
\((0,1)\cup(9,+\infty)\)
\end{answer}

\begin{solution}
若方程有实数根，则左边的非负量 \(|x^2+3x|\) 等于 \(a|x-1|\)，所以必须有 \(a>0\)；\(a=0\) 时只有 \(x=-3,0\) 两个根，\(a<0\) 时无根. 对 \(a>0\)，两边平方不改变等价性，得
\[
(x^2+3x)^2=a^2(x-1)^2,
\]
即
\[
x^2+3x=a(x-1)\quad\text{或}\quad x^2+3x=-a(x-1).
\]
这给出两个二次方程
\[
x^2+(3-a)x+a=0,
\]
\[
x^2+(3+a)x-a=0.
\]
第二个方程的判别式为 \((3+a)^2+4a=(a+1)(a+9)>0\)，恒有两个不相同的实根. 第一个方程的判别式为
\[
(3-a)^2-4a=(a-1)(a-9).
\]
要使它也有两个不相同的实根，需 \(0<a<1\) 或 \(a>9\). 两个二次方程不可能有公共根，因为公共根将同时满足 \(a(x-1)=0\) 和 \(x(x+3)=0\)，而 \(a>0\) 时这不可能. 因而恰有四个互异实根的充要条件为
\[
a\in(0,1)\cup(9,+\infty).
\]
\end{solution}

\section{解答题}

\begin{problem}
已知函数 \(f(x) = \cos x \cdot \sin\left(x + \frac{\pi}{3}\right) - \sqrt{3}\cos^2 x + \frac{\sqrt{3}}{4}\)，\(x \in \R\).

\begin{enumerate}
\item 求 \(f(x)\) 的最小正周期；
\item 求 \( f(x) \) 在闭区间 \(\left[-\frac{\pi}{4}, \frac{\pi}{4}\right]\) 上的最大值和最小值.
\end{enumerate}
\end{problem}

\begin{solution}
由和角公式及二倍角公式，
\[
\begin{gathered}
f(x)=\frac12\sin x\cos x-\frac{\sqrt3}{2}\cos^2x+\frac{\sqrt3}{4}\\
=\frac14\sin2x-\frac{\sqrt3}{4}\cos2x
=\frac12\sin\left(2x-\frac{\pi}{3}\right).
\end{gathered}
\]
因此 \(f(x)\) 的最小正周期为 \(\pi\).

当 \(x\in\left[-\frac{\pi}{4},\frac{\pi}{4}\right]\) 时，令 \(u=2x-\frac{\pi}{3}\)，则 \(u\in\left[-\frac{\pi}{2},\frac{\pi}{6}\right]\). 在此区间上 \(\sin u\) 单调递增，所以当 \(u=-\frac{\pi}{2}\)，即 \(x=-\frac{\pi}{4}\) 时，\(f(x)\) 取得最小值 \(-\frac12\)；当 \(u=\frac{\pi}{6}\)，即 \(x=\frac{\pi}{4}\) 时，\(f(x)\) 取得最大值 \(\frac14\).
\end{solution}

\begin{problem}
某大学志愿者协会有 \(6\) 名男同学，\(4\) 名女同学. 在这 \(10\) 名同学中，\(3\) 名同学来自数学学院，其余 \(7\) 名同学来自物理，化学等其他互不相同的七个学院. 现从这 \(10\) 名同学中随机选取 \(3\) 名同学，到希望小学进行支教活动（每位同学被选到的可能性相同）.

\begin{enumerate}
\item 求选出的 \(3\) 名同学是来自互不相同学院的概率；
\item 设 \(X\) 为选出的 \(3\) 名同学中女同学的人数，求随机变量 \(X\) 的分布列和数学期望.
\end{enumerate}
\end{problem}

\begin{solution}
从 \(10\) 名同学中任选 \(3\) 名，共有 \(\mathrm C_{10}^{3}\) 种等可能结果. 选出的同学来自互不相同学院，当且仅当选出的数学学院同学至多有 \(1\) 名. 因而满足条件的选法有 \(\mathrm C_{7}^{3}+\mathrm C_{3}^{1}\mathrm C_{7}^{2}\) 种，所以所求概率为
\[
P=\frac{\mathrm C_{7}^{3}+\mathrm C_{3}^{1}\mathrm C_{7}^{2}}{\mathrm C_{10}^{3}}=\frac{35+63}{120}=\frac{49}{60}.
\]

设选出的 \(3\) 名同学中女同学人数为 \(X\). 对于 \(k=0,1,2,3\)，从 \(4\) 名女同学中选出 \(k\) 名、从 \(6\) 名男同学中选出 \(3-k\) 名，因此
\[
P(X=k)=\frac{\mathrm C_{4}^{k}\mathrm C_{6}^{3-k}}{\mathrm C_{10}^{3}}.
\]
所以 \(X\) 的分布列为
\begin{center}
\begin{tabular}{c|cccc}
\hline
\(X\) & \(0\) & \(1\) & \(2\) & \(3\) \\
\hline
\(P(X=k)\) & \(\frac16\) & \(\frac12\) & \(\frac3{10}\) & \(\frac1{30}\) \\
\hline
\end{tabular}
\end{center}
并且
\[
E(X)=0\cdot\frac16+1\cdot\frac12+2\cdot\frac3{10}+3\cdot\frac1{30}=\frac65.
\]
\end{solution}

\begin{problem}
如图，在四棱锥 \(P - ABCD\) 中，\(PA \perp\) 底面 \(ABCD\)，\(AD \perp AB\)，\(AB \parallel DC\)，\(AD = DC = AP = 2\)，\(AB = 1\)，点 \(E\) 为棱 \(PC\) 的中点.

\begin{enumerate}
\item 证明：\(BE \perp DC\)；
\item 求直线 \(BE\) 与平面 \(PBD\) 所成角的正弦值；
\item 若 \(F\) 为棱 \(PC\) 上一点，满足 \(BF \perp AC\)，求二面角 \(F - AB - P\) 的余弦值.
\end{enumerate}
\bitmapfigure[max width=0.42\linewidth]{img/2014/tianjin_science/q17_fig1.png}
\end{problem}

\begin{solution}
建立空间直角坐标系，取 \(A\) 为原点，\(AB\)，\(AD\)，\(AP\) 分别为 \(x\) 轴，\(y\) 轴，\(z\) 轴正方向. 由题设可取 \(A(0,0,0)\)，\(B(1,0,0)\)，\(D(0,2,0)\)，\(C(2,2,0)\)，\(P(0,0,2)\).

\begin{enumerate}
\item 因为 \(E\) 是 \(PC\) 的中点，所以 \(E(1,1,1)\)，从而 \(\overrightarrow{BE}=(0,1,1)\)，\(\overrightarrow{DC}=(2,0,0)\). 于是 \(\overrightarrow{BE}\cdot\overrightarrow{DC}=0\)，故 \(BE\perp DC\).

\item 设直线 \(BE\) 与平面 \(PBD\) 所成的角为 \(\alpha\). 取平面 \(PBD\) 内的两个方向向量 \(\overrightarrow{BD}=(-1,2,0)\)，\(\overrightarrow{BP}=(-1,0,2)\). 于是平面 \(PBD\) 的一个法向量为
\[
\bs n=\overrightarrow{BD}\times\overrightarrow{BP}=(4,2,2).
\]
直线 \(BE\) 的方向向量为 \(\bs v=(0,1,1)\)，因此
\[
\sin\alpha=\frac{|\bs v\cdot\bs n|}{|\bs v|\,|\bs n|}=\frac{4}{\sqrt2\sqrt{24}}=\frac{\sqrt3}{3}.
\]

\item 设 \(F\) 为棱 \(PC\) 上一点. 令 \(F=P+t(C-P)\)，其中 \(0\leqslant t\leqslant1\)，则 \(F(2t,2t,2-2t)\)，\(\overrightarrow{BF}=(2t-1,2t,2-2t)\). 又 \(\overrightarrow{AC}=(2,2,0)\)，由 \(BF\perp AC\) 得
\[
\overrightarrow{BF}\cdot\overrightarrow{AC}=2(2t-1)+4t=0.
\]
解得 \(t=\frac14\)，所以 \(F\left(\frac12,\frac12,\frac32\right)\).

平面 \(PAB\) 的一个法向量为 \(\bs n_1=(0,1,0)\)，平面 \(FAB\) 的一个法向量为
\[
\bs n_2=\overrightarrow{AB}\times\overrightarrow{AF}=(0,-3,1).
\]
设二面角 \(F-AB-P\) 为 \(\beta\)，则
\[
\cos\beta=\frac{|\bs n_1\cdot\bs n_2|}{|\bs n_1|\,|\bs n_2|}=\frac{3}{\sqrt{10}}=\frac{3\sqrt{10}}{10}.
\]
\end{enumerate}
\end{solution}

\begin{problem}
设椭圆 \(\frac{x^2}{a^2}+\frac{y^2}{b^2}=1\)（\(a>b>0\)）的左、右焦点分别为 \(F_1\)，\(F_2\)，右顶点为 \(A\)，上顶点为 \(B\)，已知 \(|AB|=\frac{\sqrt3}{2}|F_1F_2|\).

\begin{enumerate}
\item 求椭圆的离心率；
\item 设 \(P\) 为椭圆上异于其顶点的一点，以线段 \(PB\) 为直径的圆经过点 \(F_1\)，经过原点 \(O\) 的直线 \(l\) 与该圆相切. 求直线 \(l\) 的斜率.
\end{enumerate}
\end{problem}

\begin{solution}
设椭圆的半焦距为 \(c\)，则 \(c^2=a^2-b^2\)，\(|AB|=\sqrt{a^2+b^2}\)，\(|F_1F_2|=2c\). 由题设，
\[
a^2+b^2=3c^2=3(a^2-b^2),
\]
所以 \(a^2=2b^2\)，从而 \(c^2=b^2\)，椭圆的离心率为
\[
e=\frac ca=\frac{\sqrt2}{2}.
\]

为求直线 \(l\) 的斜率，取 \(c>0\)，则 \(a^2=2c^2\)，\(b^2=c^2\)，可设 \(F_1(-c,0)\)，\(B(0,c)\)，\(P(x,y)\). 由于以 \(PB\) 为直径的圆经过 \(F_1\)，所以 \(\angle PF_1B=90^\circ\)，即
\[
\overrightarrow{F_1P}\cdot\overrightarrow{F_1B}=0.
\]
代入坐标得 \((x+c,y)\cdot(c,c)=0\)，即 \(x+y+c=0\). 又因为 \(P\) 在椭圆上，满足
\[
\frac{x^2}{2c^2}+\frac{y^2}{c^2}=1.
\]
令 \(X=x/c\)，\(Y=y/c\)，则 \(Y=-X-1\). 代入椭圆方程得
\[
\frac{X^2}{2}+(X+1)^2=1,
\]
所以 \(X\left(\frac32X+2\right)=0\). 当 \(X=0\) 时，\(P=(0,-c)\) 是椭圆的顶点，不合题意；故 \(X=-\frac43\)，从而
\[
P\left(-\frac{4c}{3},\frac c3\right).
\]

设以 \(PB\) 为直径的圆的圆心为 \(G\)，半径为 \(r\). 则 \(G\left(-\frac{2c}{3},\frac{2c}{3}\right)\)，且 \(r=\frac{|PB|}{2}=\frac{\sqrt5c}{3}\). 直线 \(l\) 经过原点，设其斜率为 \(k\)，则方程为 \(y=kx\). 该直线与圆相切，当且仅当圆心到直线的距离等于半径，即
\[
\frac{2c|k+1|}{3\sqrt{k^2+1}}=\frac{\sqrt5c}{3}.
\]
平方并化简，得
\[
4(k+1)^2=5(k^2+1),\qquad k^2-8k+1=0.
\]
所以直线 \(l\) 的斜率为
\[
k=4\pm\sqrt{15}.
\]
\end{solution}

\begin{problem}
已知 \(q\) 和 \(n\) 均为给定的大于 \(1\) 的自然数. 设集合 \(M = \{0, 1, 2, \dots, q - 1\}\)，集合 \(A = \{x \mid x = x_1 + x_2q + \dots + x_nq^{n-1}, x_i \in M, i = 1, 2, \dots, n\}\).

\begin{enumerate}
\item 当 \(q = 2\)，\(n = 3\) 时，用列举法表示集合 \(A\)；
\item 设 \(s\)，\(t \in A\)，\(s = a_1 + a_2q + \cdots + a_nq^{n-1}\)，\(t = b_1 + b_2q + \cdots + b_nq^{n-1}\)，其中 \(a_i\)，\(b_i \in M\)，\(i = 1,2,\cdots,n\). 证明：若 \(a_n < b_n\)，则 \(s < t\).
\end{enumerate}
\end{problem}

\begin{solution}
\begin{enumerate}
\item 当 \(q=2\)，\(n=3\) 时，\(x_1,x_2,x_3\) 均可取 \(0\) 或 \(1\)，所以
\[
A=\{x_1+2x_2+4x_3\mid x_1,x_2,x_3\in\{0,1\}\}=\{0,1,2,3,4,5,6,7\}.
\]

\item 由 \(a_n<b_n\) 且 \(a_n,b_n\in M\)，有 \(b_n-a_n\geqslant1\). 于是
\[
\begin{gathered}
t-s=(b_n-a_n)q^{n-1}+\sum_{i=1}^{n-1}(b_i-a_i)q^{i-1}\\
\geqslant q^{n-1}-(q-1)\sum_{i=1}^{n-1}q^{i-1}=q^{n-1}-(q^{n-1}-1)=1>0.
\end{gathered}
\]
因此 \(t-s>0\)，即 \(s<t\).
\end{enumerate}
\end{solution}

\begin{problem}
已知函数 \(f(x)=x-a\e^{x}\)（\(a\in\R\)，\(x\in\R\)），已知函数 \(y=f(x)\) 有两个零点 \(x_1\)，\(x_2\)，且 \(x_1<x_2\).

\begin{enumerate}
\item 求 \(a\) 的取值范围；
\item 证明 \(\frac{x_2}{x_1}\) 随着 \(a\) 的减小而增大；
\item 证明 \(x_1+x_2\) 随着 \(a\) 的减小而增大.
\end{enumerate}
\end{problem}

\begin{solution}
零点满足 \(x-a\e^x=0\)，即 \(a=x\e^{-x}\). 若 \(a<0\)，则 \(f'(x)=1-a\e^x>0\)，函数严格递增，只能有一个零点；若 \(a=0\)，则 \(f(x)=x\)，也只有一个零点. 因而要有两个零点，必须有 \(a>0\)，此时零点均为正数.

令 \(g(x)=x\e^{-x}\)（\(x>0\)），则 \(g'(x)=\e^{-x}(1-x)\). 所以 \(g(x)\) 在 \((0,1)\) 上递增，在 \((1,+\infty)\) 上递减，且 \(g(1)=\frac1{\e}\). 因此直线 \(y=a\) 与曲线 \(y=g(x)\) 有两个交点的充要条件是
\[
0<a<\frac1{\e}.
\]
此时 \(0<x_1<1<x_2\).

由 \(x_i=a\e^{x_i}\) 关于 \(a\) 求导，得
\[
(1-x_i)\frac{\mathrm d x_i}{\mathrm d a}=\e^{x_i}=\frac{x_i}{a},\qquad \frac{\mathrm d x_i}{\mathrm d a}=\frac{x_i}{a(1-x_i)}.
\]
于是
\[
\frac{\mathrm d}{\mathrm d a}\ln\frac{x_2}{x_1}=\frac{1}{a(1-x_2)}-\frac{1}{a(1-x_1)}=\frac{x_2-x_1}{a(1-x_1)(1-x_2)}<0.
\]
所以 \(x_2/x_1\) 随 \(a\) 的增大而减小，等价地，随 \(a\) 的减小而增大.

再设 \(S=x_1+x_2\)，则
\[
\frac{\mathrm d S}{\mathrm d a}=\frac{x_1}{a(1-x_1)}+\frac{x_2}{a(1-x_2)}=\frac{x_1+x_2-2x_1x_2}{a(1-x_1)(1-x_2)}.
\]
由 \(x_1\e^{-x_1}=x_2\e^{-x_2}\)，可得 \(\ln\frac{x_2}{x_1}=x_2-x_1\). 令 \(u=x_2/x_1>1\)，则
\[
x_1=\frac{\ln u}{u-1},\qquad x_2=\frac{u\ln u}{u-1}.
\]
又令 \(\varphi(u)=\frac{u-u^{-1}}{2}-\ln u\)，则 \(\varphi(1)=0\)，且
\[
\varphi'(u)=\frac{(u-1)^2}{2u^2}>0\quad(u>1).
\]
故 \(u-u^{-1}>2\ln u\)，从而
\[
x_1+x_2-2x_1x_2=\frac{\ln u}{(u-1)^2}\left(u^2-1-2u\ln u\right)>0.
\]
由于 \(a>0\)，\(1-x_1>0\)，\(1-x_2<0\)，所以 \(\frac{\mathrm d S}{\mathrm d a}<0\). 因此 \(x_1+x_2\) 随 \(a\) 的减小而增大.
\end{solution}
